%% =====================================================================
%%  Corrigés — Unit 2 « Where are you from? »
%% =====================================================================
\corrigesUnite{2}{Where are you from?}

\begin{corrige}{1}
\begin{reponses}[4]
\item \rep{Spanish}
\item \rep{Chinese}
\item \rep{Italian}
\item \rep{Brazilian}
\item \rep{Mexico}
\item \rep{Ireland}
\item \rep{Canada}
\item \rep{Germany}
\end{reponses}
\retiens{pays et nationalités prennent toujours une majuscule : \en{Spain}, \en{Spanish}.}
\end{corrige}

\begin{corrige}{2}
\begin{reponses}[2]
\item \rep{I'm not Italian.}
\item \rep{Jack isn't British.}
\item \rep{We aren't from Paris.}
\item \rep{Emma isn't American.}
\item \rep{They aren't Japanese.}
\item \rep{It isn't a big city.}
\end{reponses}
\pourquoi{on peut aussi contracter le sujet et le verbe : \en{We're not from Paris}, \en{They're
not Japanese}, \en{It's not a big city}. Les deux formes sont correctes. Avec \en{I}, une seule
possibilité : \en{I'm not}.}
\end{corrige}

\begin{corrige}{3}
\begin{reponses}[2]
\item \rep{Is Léa French?}
\item \rep{Are you from Brooklyn?}
\item \rep{Is Emma a beginner?}
\item \rep{Are they Canadian?}
\item \rep{Is Mrs. Lopez Mexican?}
\item \rep{Is Lyon a big city?}
\end{reponses}
\pourquoi{le verbe \en{be} passe devant le sujet, qu'il s'agisse d'un pronom (\en{Are you…?}) ou
d'un nom (\en{Is Léa…?}). Le reste de la phrase ne bouge pas.}
\piege{oublier le point d'interrogation, ou garder l'ordre de l'affirmation : \en{Léa is French?}
n'est pas une vraie question écrite.}
\end{corrige}

\begin{corrige}{4}
\begin{reponses}[2]
\item \rep{Yes, she is.}
\item \rep{No, he isn't.}
\item \rep{Yes, she is.}
\item \rep{No, they aren't.}
\item \rep{Yes, she is.}
\item \rep{Yes, it is.}
\end{reponses}
\pourquoi{la réponse courte reprend le verbe de la question avec un pronom : \en{Is Léa…?} →
\en{she is}. Jack est américain (pas britannique) ; Jack et Emma viennent de Brooklyn et de
Londres ; Mrs.\ Lopez est américaine ; Léa dit de Lyon : \en{It's a beautiful city!}}
\piege{contracter une réponse courte affirmative : on dit \en{Yes, she is}, jamais \en{Yes, she's}.}
\end{corrige}

\begin{corrige}{5}
\begin{reponses}[3]
\item \rep{is}
\item \rep{am}
\item \rep{isn't}
\item \rep{are}
\item \rep{aren't}
\item \rep{is}
\end{reponses}
\pourquoi{\en{my best friend} et \en{it} sont au singulier (\en{is}, \en{isn't}) ; \en{my parents}
et \en{they} au pluriel (\en{are}, \en{aren't}). Le sens guide le choix entre affirmation et
négation : l'amie d'Emma est irlandaise, donc elle n'est pas britannique.}
\end{corrige}

\begin{corrige}{6}
\begin{reponses}[1]
\item \rep{Léa is French.}
\item \rep{Jack speaks English and a little Spanish.}
\item \rep{Emma is from London, in England.}
\item \rep{I'm not Italian.}
\end{reponses}
\retiens{majuscule en anglais pour : le début de la phrase, \en{I}, les prénoms, les villes, les
pays, les nationalités et les langues.}
\end{corrige}

\begin{corrige}{7}
\begin{reponses}[4]
\item \rep{ja-PAN}
\item \rep{ja-pa-NESE}
\item \rep{IT-a-ly}
\item \rep{i-TAL-ian}
\item \rep{CAN-a-da}
\item \rep{ca-NA-di-an}
\item \rep{a-MER-i-ca}
\item \rep{bra-ZIL}
\end{reponses}
\pourquoi{la syllabe accentuée est en capitales. Comparez \en{Japan} et \en{Japanese}, \en{Italy}
et \en{Italian}, \en{Canada} et \en{Canadian} : l'accent se déplace quand on passe du pays à la
nationalité. Texte de la piste page~\pageref{trans-2.10}.}
\end{corrige}

\begin{corrige}{8}
\begin{reponses}[2]
\item \rep{Brazilian}
\item \rep{Rio de Janeiro}
\item \rep{Italian}
\item \rep{Rome}
\item \rep{Japanese}
\item \rep{Tokyo}
\item \rep{Irish}
\item \rep{Dublin}
\end{reponses}
\pourquoi{Yuki ne donne pas sa nationalité, mais elle vient de Tokyo, au Japon : elle est
\en{Japanese}. Liam précise : \en{I'm not English, I'm Irish!}}
\end{corrige}

\begin{corrige}{9}
\begin{anglais}\otTableau
\begin{tabular}{@{}l@{\hspace{6mm}}l@{}}
\otNomDit{sam} & Hi, Léa! \rep{Are} you OK?\\
\otNomDit{lea} & Yes, \rep{I'm} fine, thanks.\\
\otNomDit{sam} & Are you \rep{from} Paris?\\
\otNomDit{lea} & No, I'm \rep{not}. I'm from Lyon.\\
\otNomDit{sam} & Is Lyon a \rep{big} city?\\
\otNomDit{lea} & Yes, \rep{it} is. And you? Where are you from?\\
\otNomDit{sam} & I'm \rep{from} here, from Brooklyn!\\
\otNomDit{lea} & \rep{Really}? Great!\\
\end{tabular}
\end{anglais}
\retiens{\en{Are you OK?} = « Ça va ? » (« tu vas bien ? »). On répond \en{Yes, I'm fine, thanks.}}
\end{corrige}

\begin{corrige}{10}
\begin{reponses}[1]
\item \rep{I'm from France.}
\item \rep{Are you American? — No, I'm not.}
\item \rep{She isn't Spanish, she's Mexican.}
\item \rep{Where are you from?}
\item \rep{I speak French and a little English.}
\item \rep{They aren't from London.}
\item \rep{London is a big city.}
\item \rep{Is Emma British? — Yes, she is.}
\end{reponses}
\pourquoi{phrase 3 : on peut aussi écrire \en{She's not Spanish}. Phrase 6 : \en{They're not
from London} est également correct.}
\piege{traduire « je viens de » par \en{I come of}. On dit simplement \en{I'm from}.}
\end{corrige}

\begin{corrige}{11}
\begin{reponses}[3][french]
\item \rep{F}
\item \rep{V}
\item \rep{F}
\item \rep{V}
\item \rep{V}
\item \rep{F}
\end{reponses}
\pourquoi{Léa vient de Lyon et elle est maintenant à New York (\en{Now I'm in New York}) ; Emma
est britannique : \en{Emma isn't American}.}
\end{corrige}

\begin{corrige}{12}
Réponses modèles (adaptez-les à votre situation) :
\begin{reponses}[1]
\item \rep{I'm from France.} \textit{ou} \rep{I'm from Belgium.}
\item \rep{I'm from Lyon.} \textit{ou} \rep{I'm from Marseille.}
\item \rep{No, I'm not. I'm French.}
\item \rep{Yes, I am.}
\item \rep{Thank you!} \textit{ou} \rep{Thanks, Emma!}
\end{reponses}
\retiens{à la question 3, la réponse courte (\en{No, I'm not.}) suivie de l'information juste
(\en{I'm French.}) est la réponse la plus naturelle en anglais.}
\end{corrige}

\begin{corrigetest}
\begin{reponses}[5][french]
\item \rep{a}
\item \rep{b}
\item \rep{a}
\item \rep{b}
\item \rep{b}
\item \rep{c}
\item \rep{b}
\item \rep{c}
\item \rep{a}
\item \rep{a}
\end{reponses}
\pourquoi{question 4 : on inverse le verbe et le sujet (\en{Are you…?}). Question 7 : les langues
prennent une majuscule. Question 9 : \en{the UK} s'écrit avec \en{the}, car \en{Kingdom} est un
nom commun.}
\end{corrigetest}
